% ID: SEL-D0083-OB-E01
% ESTADO: cotejada. Segunda pasada visual y borrador compilado revisados.
% Fuente DOC: Selectividad/sel_2012_matematicas_aplicadas/5 MATEMÁT. APLIC. CC. SS/EXAMEN 3-ANDALUCÍA.doc
% SHA-256 original: c8688a19515dfeb348e8f8817d3f4a6c0a967723c270f2ec42e7f365d001e555
% Página convertida: 2; opción B; ejercicio 1.

Una fábrica produce dos tipos de productos, A y B, que distribuye a tres clientes. En el mes de enero el primer cliente compró 9 unidades de A y 5 de B, el segundo cliente 3 de A y 7 de B, y el tercer cliente 4 de A y 6 de B.

En el mes de febrero el primer cliente y el segundo duplicaron las compras del mes anterior, y el tercer cliente compró de cada producto una unidad más de las que compró en enero. En marzo el primer cliente no compró nada, y el segundo y el tercero compraron lo mismo que en febrero.
\begin{enumerate}[label=\alph*)]
\item \textbf{(0.75 puntos)} Para cada mes construya la matriz de dimensión $3\mathrm{x}2$ correspondiente a las compras de ese mes.
\item \textbf{(0.5 puntos)} Calcule la matriz de compras del trimestre.
\item \textbf{(1.25 puntos)} Si los precios de los productos A y B son, respectivamente, 80 y 100 euros, calcule lo que factura la fábrica en el primer trimestre, por cada cliente y en total.
\end{enumerate}
